\documentclass[exercice,a4]{thedocument}%
\usepackage{tabularx}%
\startdatakeys
    \source{25GENMATPO3}
    \BO{2026}
    \cycle{4}
    \competencesDuSocle{CA,RA}
    \objectifsApprentissage{C4A1aa,C4A2ad,C4B2cb,C4D2aa,C4D3ac}
    \motsCles{pourcentage,tableur,probabilité,moyenne,médiane,PGCD,facteurs premiers}
\enddatakeys
\begin{document}%
\medskip

Une étude montre qu’un nombre important de nos vêtements reste dans les placards et n’est pas utilisé. Six amis décident de donner chacun une partie de leurs vêtements à une association lors d’une journée de collecte organisée dans leur village.

\medskip

\begin{enumerate}
\item Inès compte $20$ tee-shirts dans son placard et souhaite en donner $70\,\%$ à l’association. Montrer qu’elle va en donner 14.
\item La feuille de calcul ci-dessous indique le nombre de tee-shirts que chacun des amis souhaite donner à cette association.

\begin{center}
\begin{tabularx}{\linewidth}{|c|m{3cm}|*{7}{>{\centering \arraybackslash}X|}}\hline
&A&B&C&D&E&F&G&H\\ \hline
1&&Inès&Sylvain&Sabrina&Marco&Yuna&Axel&Total\\ \hline
2&Nombre de tee-shirts à donner&14&6&9&11&12&8&\\ \hline
\end{tabularx}
\end{center}

	\begin{enumerate}
		\item Quelle formule peut-on saisir dans la cellule H2 pour obtenir le nombre total de tee-shirts à donner ?
		\item Les amis réunissent tous les tee-shirts qu’ils vont donner.
		
Calculer la probabilité qu’un tee-shirt pris au hasard appartienne à Yuna.
		\item Calculer le nombre moyen de tee-shirts donnés par chacun des amis.
		\item Quelle est la médiane du nombre de tee-shirts donnés ?
	\end{enumerate}
\item A la fin de la journée de collecte des vêtements, l’association a récolté un total de $168$ tee-shirts et $63$ pantalons. Cette association souhaite réaliser des lots identiques contenant chacun le même nombre de pantalons et le même nombre de tee-shirts en utilisant tous les vêtements donnés.
	\begin{enumerate}
		\item Peut-elle réaliser 4 lots ? Peut-elle réaliser 3 lots ?\index{P{}. G. C. D.}
		\item Déterminer la décomposition en produit de facteurs premiers de $168$ et de $63$.
		\item Quel nombre maximum de lots pourra-t-elle réaliser ?
	\end{enumerate}
\end{enumerate}

\newpage

\begin{center}
\textbf{\large ANNEXE à rendre avec votre copie}
\end{center}
\showthesource
\end{document}
